%! Introduction to Polynomial Functions — Unit 2, Lesson 2.0 — Exit Ticket
\documentclass[10pt]{article}
\usepackage{precalculus-article}
\usepackage{precalculus-boxes}
\newcommand{\TallMath}[1]{$\displaystyle #1\rule[-1.4em]{0pt}{3.2em}$}

\begin{document}

\pageheader{Unit 2, Lesson 2.0}{Exit Ticket}

\vspace{0.12in}

\begin{enumerate}[itemsep=12pt]

\item Consider the polynomial $p(x) = -4x^3 + 7x^2 - x + 9$.

\vspace{0.06in}
\begin{enumerate}[label=(\alph*), itemsep=6pt]
  \item Degree: \blank{2.0cm} \qquad Leading coefficient: \blank{2.0cm}
        \qquad Constant term: \blank{2.0cm}
  \item By degree, $p$ belongs to the family called \blank{3.0cm}.
\end{enumerate}

\item Is each expression a polynomial? Circle one. Then explain the \textbf{No} in
      one sentence.

\begin{enumerate}[label=(\alph*), itemsep=6pt]
  \item $g(x) = 6x^4 - 2x + 1$ \quad \textbf{Yes} / \textbf{No}
  \item $h(x) = 3x^2 + \dfrac{5}{x}$ \quad \textbf{Yes} / \textbf{No}
\end{enumerate}

\vspace{0.06in}
Reason for the one that is \textbf{not} a polynomial:

\vspace{0.06in}
\writeline

\item Write $5 - 2x + 4x^2$ in standard form.

\vspace{0.1in}
Standard form: \blank{5cm}

\end{enumerate}

\end{document}
